%% LyX 2.4.4 created this file.  For more info, see https://www.lyx.org/.
%% Do not edit unless you really know what you are doing.
\documentclass[english]{article}
\usepackage[T1]{fontenc}
\usepackage[utf8]{inputenc}
\setcounter{secnumdepth}{-2}
\setlength{\parindent}{0bp}
\usepackage{amsmath}
\usepackage{amsthm}
\usepackage{amssymb}

\makeatletter
%%%%%%%%%%%%%%%%%%%%%%%%%%%%%% Textclass specific LaTeX commands.
\theoremstyle{definition}
\newtheorem*{problem*}{\protect\problemname}

\@ifundefined{date}{}{\date{}}
%%%%%%%%%%%%%%%%%%%%%%%%%%%%%% User specified LaTeX commands.
\usepackage{graphicx}
\usepackage{geometry}
\newcommand{\limi}{\underset{n\rightarrow\infty}{\lim}}
\newcommand{\limxz}{\underset{x\rightarrow x_{0}}{\lim}}
\newcommand{\limfxz}{\underset{x\rightarrow x_{0}}{\lim}f(x)}
\newcommand{\limfxm}{\underset{x\rightarrow x_{0}^{-}}{\lim}f(x)}
\newcommand{\limfxp}{\underset{x\rightarrow x_{0}^{+}}{\lim}f(x)}
\newcommand{\limxm}{\underset{x\rightarrow x_{0}^{-}}{\lim}}
\newcommand{\limxp}{\underset{x\rightarrow x_{0}^{+}}{\lim}}
\newcommand{\sqrm}{M_{m\times m}(\mathbb{F})}
\newcommand{\seqa}{(a_{n})_{n=1}^{\infty}}
\newcommand{\seqx}{(x_{n})_{n=1}^{\infty}}
\newcommand{\funcdr}{f:D\rightarrow\mathbb{R}}
\newcommand{\der}{\limxz\frac{f(x)-f(x_{0})}{x-x_{0}}}
\usepackage[all]{pdfprivacy}

\makeatother

\usepackage{babel}
\providecommand{\problemname}{Problem}
\begin{document}
\newgeometry{top=2cm, bottom=2cm, left=2cm, right=2cm}
\title{{\Large Semester 2026A, Exam A, Q8\vspace{-2em}}}

\maketitle
\begin{problem*}
Let $A,B\in M_{m\times n}(\mathbb{R})$ be matrices
and let $b\in\mathbb{R}^{m}$, such that the system of linear equations
$Ax=b$ has a unique solution. Which of the following claims are true?
If a claim is true, prove it; if not, give a counterexample. 
\begin{enumerate}
\item If there exists an invertible matrix $P\in M_{m\times m}(\mathbb{R})$
such that $B=PA$, then the system of linear equations $Bx=b$ has
a unique solution. 
\item If there exists an invertible matrix $P\in M_{n\times n}(\mathbb{R})$
such that $B=AP$, then the system of linear equations $Bx=b$ has
a unique solution. 
\end{enumerate}
\end{problem*}

\medskip{}

\rule[0.5ex]{1\columnwidth}{1pt}

\medskip{}

\textbf{Analyzing the question:}

\medskip{}

By the guidance of a gifted mathematical monk I happened to come across
upon my ways, we shall approach this question from the viewpoint of
functions, as an important essence of linear algebra is that every
matrix can be represented as a linear transformation.

In particular, the matrix $A$ defines a linear transformation $T_{A}:\mathbb{R}^{n}\rightarrow\mathbb{R}^{m}$
that operates according to the rule $T_{A}(x)=Ax$. 

The given data devilishly appears as one fact but is actually two:
$Ax=b$ has a solution, and this solution is unique. 

Its existence, translated to our language of choice, means that $b\in\text{Im}(T_{A})$,
because there exists a solution s.t. $T_{A}(x)=Ax=b$. 

For uniqueness, let $U$ denote the solution set of the homogeneous
system $Ax=0$, and recall that $U=\ker$$\left(T_{A}\right)$ by
definition. Let us recall the theorem: if $Ax=b$ has a solution and
$v$ is any solution of it, then its solution set is exactly $v+U$.

Our system has a solution, and its solution set is a single vector,
say $\{v\}$. 

Applying the theorem, $v+U=\{v\}$, and hence $U=\{0\}$.

Let us recall that a linear transformation is injective if and only
if its kernel is zero, which means $T_{A}$ is injective.

We shall also recall that if $G\in M_{k\times k}(\mathbb{R})$ is
invertible, then $T_{G}:\mathbb{R}^{k}\rightarrow\mathbb{R}^{k}$
is bijective.

With these arrows in our quiver we shall bravely approach this question.

\medskip{}

\textbf{Part 1:}

\medskip{}

Part 1 asks us (and quite impolitely so) to study the linear transformation
$T_{P}\circ T_{A}$, since $B=PA$ gives $T_{B}=T_{P}\circ T_{A}$.
Do we know that there exists exactly one $x\in\mathbb{R}^{n}$ s.t.
$T_{P}\left(T_{A}\left(x\right)\right)=b$? To quote Monty Python's
argument skit, not necessarily ($T_{A}$ can choose to satisfy this
in his spare time). 

Note first what cannot go wrong. 

$T_{A}$ is injective, and $T_{P}$ is bijective, so $T_{B}=T_{P}\circ T_{A}$
is a composition of injective maps and is therefore injective. 

Hence $Bx=b$ has at most one solution, one way or another.

Only existence can fail, and fail it does.

Let us see where the counter-example can live.

If $m=n$, then $A$ is injective and square, hence invertible, so
$B=PA$ is invertible as well and $Bx=b$ has a unique solution for
every b. 

Moreover, because $T_{A}$ is injective, we must choose $n\le m$
- lest this attribute will not be satisfied.

Therefore, we must choose $m>n$.

Let us choose $A=b=\begin{pmatrix}2\\
1
\end{pmatrix}$ and $P=\begin{pmatrix}0 & 1\\
1 & 0
\end{pmatrix}$.

First, we verify that the counter-example fulfills the requirements.

Here $A\in M_{2\times1}(\mathbb{R})$ and $b\in\mathbb{R}^{2}$, and
$Ax=b$ means 

\[
\begin{pmatrix}2\\
1
\end{pmatrix}x=\begin{pmatrix}2\\
1
\end{pmatrix}\therefore x=1
\]

so $x=1$ is indeed the unique solution.

$P$ is invertible: computing $P^{-1}$ works 

\[
\begin{pmatrix}\begin{array}{cc|cc}
0 & 1 & 1 & 0\\
1 & 0 & 0 & 1
\end{array}\end{pmatrix}\overset{R_{1}\Leftrightarrow R_{2}}{\Longrightarrow}\begin{pmatrix}\begin{array}{cc|cc}
1 & 0 & 0 & 1\\
0 & 1 & 1 & 0
\end{array}\end{pmatrix}
\]

so $P^{-1}=P$, which we confirm by

\[
\begin{pmatrix}0 & 1\\
1 & 0
\end{pmatrix}\begin{pmatrix}0 & 1\\
1 & 0
\end{pmatrix}=\begin{pmatrix}\left(0\cdot0+1\cdot1\right) & \left(0\cdot1+1\cdot0\right)\\
\left(1\cdot0+0\cdot1\right) & \left(1\cdot1+0\cdot0\right)
\end{pmatrix}=\begin{pmatrix}1 & 0\\
0 & 1
\end{pmatrix}=I_{2}
\]
However, the conclusion fails: $T_{P}\left(T_{A}\left(x\right)\right)=b$
reads

\[
T_{P}\left(T_{A}\left(x\right)\right)=T_{P}\left(\begin{pmatrix}2\\
1
\end{pmatrix}x\right)=T_{P}\left(\begin{pmatrix}2x\\
x
\end{pmatrix}\right)=\begin{pmatrix}0 & 1\\
1 & 0
\end{pmatrix}\begin{pmatrix}2x\\
x
\end{pmatrix}=\begin{pmatrix}0\cdot2x+1\cdot x\\
1\cdot2x+0\cdot x
\end{pmatrix}=\begin{pmatrix}x\\
2x
\end{pmatrix}
\]

which gives $x=2,2x=4\neq1$, thus the system $Bx=b$ has no solution
at all and our counter-example is finished - claim 1 is false.

\medskip{}

\textbf{Part 2:}

\medskip{}

Part 2 kindly asks us to study the linear transformation $T_{A}\circ T_{P}$,
since $B=AP$ gives $T_{B}=T_{A}\circ T_{P}$.

Do we know that there exists exactly one $x\in\mathbb{R}^{n}$ s.t.
$T_{A}\left(T_{P}\left(x\right)\right)=b$?

This time, we do.

We pick up our bow and shoot the first arrow: we know that there is
exactly one $v\in\mathbb{R}^{n}$ with $T_{A}(v)=b$.

We shoot the second arrow: since $P$ is invertible, $T_{P}:\mathbb{R}^{n}\rightarrow\mathbb{R}^{n}$
is bijective. 

Now let $x\in\mathbb{R}^{n}$.

\[
Bx=b\iff T_{A}\left(T_{P}\left(x\right)\right)=b\underset{\text{from the uniqueness of }v}{\iff}T_{P}\left(x\right)=v
\]

But since $T_{P}$ is bijective and $v\in\mathbb{R}^{n}$, there exists
a unique $x\in\mathbb{R}^{n}$ s.t. $T_{P}(x)=v$, i.e. $x=T_{P}^{-1}(v)$.
So this x is the one and only solution, and our dragon is slain -
claim 2 is true.
\begin{quote}
This is an exam problem I translated and solved. The original exam
was written by the Linear Algebra 1 course lecturers at HUJI.
\end{quote}

\end{document}
